Los requisitos no funcionales describen las propiedades de calidad y restricciones que debe cumplir el sistema, más allá de las funcionalidades que proporciona. Estos requisitos incluyen aspectos como el rendimiento, la seguridad, la usabilidad o la mantenibilidad del sistema.

Para la clasificación de los requisitos no funcionales se ha tomado como referencia el modelo de calidad definido en el estándar ISO/IEC 25010, que establece un conjunto de características para evaluar la calidad de los sistemas software \cite{iso25010}.

% ------------------------------------------------------------------------------
\subsection{RNF1. Usabilidad}
% ------------------------------------------------------------------------------

\vspace{-0.75em}
\subsubsection{RNF-U-01 Facilidad de uso}
\vspace{-1em}

El sistema debe proporcionar una interfaz intuitiva que permita a los usuarios realizar las tareas principales con un número reducido de interacciones.

\vspace{-0.75em}
\subsubsection{RNF-U-02 Adaptabilidad de la interfaz}
\vspace{-1em}

La interfaz del sistema debe adaptarse correctamente a distintos tamaños de pantalla, incluyendo dispositivos móviles, tabletas y ordenadores, manteniendo una experiencia de usuario consistente.

\vspace{-0.75em}
\subsubsection{RNF-U-03 Claridad de la información}
\vspace{-1em}

El sistema debe presentar la información de forma clara y estructurada, facilitando la comprensión de los datos mostrados por parte de los usuarios.

% ------------------------------------------------------------------------------
\subsection{RNF2. Eficiencia de rendimiento}
% ------------------------------------------------------------------------------

\vspace{-0.75em}
\subsubsection{RNF-P-01 Tiempo de respuesta}
\vspace{-1em}

El sistema debe responder a las solicitudes de los usuarios en un tiempo nferior a 2 segundos en condiciones normales de uso.

\vspace{-0.75em}
\subsubsection{RNF-P-02 Optimización de recursos}
\vspace{-1em}

El sistema debe optimizar la carga de recursos para garantizar un funcionamiento fluido incluso en dispositivos móviles con capacidad limitada.

% ------------------------------------------------------------------------------
\subsection{RNF3. Compatibilidad}
% ------------------------------------------------------------------------------

\vspace{-0.75em}
\subsubsection{RNF-C-01 Compatibilidad con navegadores}
\vspace{-1em}

El sistema debe ser compatible con los navegadores web modernos más utilizados.

\vspace{-0.75em}
\subsubsection{RNF-C-02 Compatibilidad con sistemas operativos móviles}
\vspace{-1em}

El sistema debe ser compatible con los sistemas operativos móviles más utilizados (android e iOS).

\vspace{-0.75em}
\subsubsection{RNF-C-03 Integración con servicios externos}
\vspace{-1em}

El sistema debe permitir la integración con servicios externos necesarios para el funcionamiento del sistema, tales como servicios de mapas o herramientas de gestión empresarial, así como automatizar la integración con dichos servicios para garantizar su correcto funcionamiento y actualización.

% ------------------------------------------------------------------------------
\subsection{RNF4. Fiabilidad}
% ------------------------------------------------------------------------------

\vspace{-0.75em}
\subsubsection{RNF-F-01 Disponibilidad del sistema}
\vspace{-1em}

El sistema debe estar disponible para los usuarios la mayor parte del tiempo, minimizando interrupciones en el servicio.

\vspace{-0.75em}
\subsubsection{RNF-F-02 Integridad de los datos}
\vspace{-1em}

El sistema debe garantizar la correcta conservación de los datos almacenados, evitando pérdidas o inconsistencias en la información.

% ------------------------------------------------------------------------------
\subsection{RNF5. Seguridad}
% ------------------------------------------------------------------------------

\vspace{-0.75em}
\subsubsection{RNF-S-01 Control de acceso}
\vspace{-1em}

El sistema debe garantizar que únicamente los usuarios autorizados puedan acceder a las funcionalidades correspondientes a su rol.

\vspace{-0.75em}
\subsubsection{RNF-S-02 Protección de datos}
\vspace{-1em}

El sistema debe garantizar la protección de los datos almacenados mediante mecanismos adecuados de autenticación y control de acceso.

% ------------------------------------------------------------------------------
\subsection{RNF6. Mantenibilidad}
% ------------------------------------------------------------------------------

\vspace{-0.75em}
\subsubsection{RNF-M-01 Modularidad del sistema}
\vspace{-1em}

El sistema debe estar organizado en módulos independientes que faciliten su mantenimiento, evolución y ampliación.

\vspace{-0.75em}
\subsubsection{RNF-M-02 Facilidad de actualización}
\vspace{-1em}

El sistema debe permitir actualizar o modificar funcionalidades sin afectar al funcionamiento global de la aplicación.

% ------------------------------------------------------------------------------
\subsection{RNF7. Portabilidad}
% ------------------------------------------------------------------------------

\vspace{-0.75em}
\subsubsection{RNF-PO-01 Acceso multiplataforma}
\vspace{-1em}

El sistema debe poder ejecutarse en distintos dispositivos mediante un navegador web sin necesidad de instalaciones específicas.

\vspace{-0.75em}
\subsubsection{RNF-PO-02 Instalación como aplicación web}
\vspace{-1em}

El sistema debe permitir su instalación como aplicación web progresiva (PWA) en dispositivos compatibles, facilitando su uso como aplicación móvil.

